\begin{table}[!b]
\centering\normalsize
\caption{Workflow-native authorization taxonomy. Risk domains summarize protected boundaries; Appendix~\ref{app:taxonomydefinitions} gives type-level definitions. Authorization mode is assigned separately.}
\label{tab:taxonomy}
\begin{tabularx}{\linewidth}{@{}>{\raggedright\arraybackslash}p{.21\linewidth}>{\raggedright\arraybackslash}p{.27\linewidth}>{\raggedright\arraybackslash}X@{}}
\toprule
\textbf{Risk domain} & \textbf{Protected boundary} & \textbf{Attack types} \\
\midrule
Data confidentiality
& Permitted use of protected data
& Access-to-Disclosure; Excessive Data Extraction; Unauthorized Data Reuse \\
\midrule
Recipient and destination
& Authorized recipients and endpoints
& Additional Destination; Audience Expansion; Cross-Boundary Redirection \\
\midrule
Authorization validity
& Valid principal, delegation, and time
& Stale or Revoked Authorization; Role-as-Authorization; Invalid Approver \\
\midrule
Approval scope and binding
& Action and object named in an approval
& Cross-Action Approval Transfer; Wrong-Object Approval \\
\midrule
Privilege and policy
& Least privilege and mandatory policy
& Privilege Escalation; Hard-Policy Bypass \\
\midrule
Operational scope
& Duration, count, and resource limits
& Unauthorized Persistence; Repeated Execution; Resource Overuse \\
\midrule
Execution bypass
& Effect-level enforcement across tools and executors
& Cross-Tool Bypass; Delegated Execution \\
\midrule
Audit and accountability
& Observable evidence and attributable principals
& Audit Suppression; Attribution Obfuscation \\
\bottomrule
\end{tabularx}
\end{table}
